%----------------------------------------------------------------------
% 文獻討討
%----------------------------------------------------------------------

\chapter{Related Work}\label{chap:related_works}

\section{Multi-Token Prediction}\label{sec:2-mtp}

Multi-Token Prediction (MTP) extends conventional next-token prediction by enabling a language model to predict multiple future tokens simultaneously. Gloeckle et al.\cite{gloeckle2024mtp} proposed an MTP framework consisting of a shared model trunk and multiple prediction heads, where each head is responsible for predicting a different future token. Their results showed that MTP can improve model performance and support self-speculative decoding for faster inference.

Samragh et al.\cite{samragh2025future} further explored introducing multi-token prediction capabilities into pretrained autoregressive language models through supervised fine-tuning. Their method incorporates additional training objectives, including consistency learning, to improve the accuracy and coherence of jointly predicted future tokens while enabling faster decoding.

\section{Speculative Decoding}\label{sec:2-speculative_decoding}

Stern et al.\cite{stern2018blockwise} proposed Blockwise Parallel Decoding, an early approach for accelerating autoregressive generation by predicting multiple future tokens in parallel. The predicted tokens are then validated by a scoring model, and the longest valid prefix is accepted at each decoding step.

Leviathan et al.\cite{leviathan2023speculative} introduced speculative decoding, in which a lightweight draft model generates multiple candidate tokens that are subsequently verified in parallel by a larger target model. This approach reduces the number of sequential decoding steps while preserving the output distribution of the target model.

Cai et al.\cite{cai2024medusa} proposed Medusa, which eliminates the need for a separate draft model by adding multiple decoding heads to the target model. These heads predict multiple future tokens in parallel, while a tree-based attention mechanism is used to construct and verify candidate continuations efficiently.

\section{Consistency Learning}\label{sec:2-consistency_learning}

Consistency learning has been widely used as a regularization strategy to improve the stability and generalization of neural networks. Liang et al.\cite{liang2021rdrop} proposed R-Drop, which minimizes the bidirectional KL divergence between the output distributions of two sub-models generated by dropout, thereby encouraging consistent predictions under different stochastic perturbations.

Ni et al.\cite{ni2024lrdrop} extended consistency regularization from the output level to intermediate representations in Transformer-based language models. Their LR-Drop method aligns hidden states, attention matrices, and output distributions between different dropout-based sub-models, demonstrating that consistency constraints can be applied at multiple representation levels.

For multi-token prediction, Samragh et al.\cite{samragh2025future} introduced auxiliary consistency objectives to improve the coherence and accuracy of jointly predicted future tokens. Their approach applies consistency learning during supervised fine-tuning, showing that explicitly constraining future-token predictions can improve multi-token generation performance.

\section{Research Gap}\label{sec:2-research_gap}

Existing studies have demonstrated the potential of multi-token prediction for improving future-token prediction and accelerating autoregressive decoding. However, standard MTP primarily optimizes each prediction head independently using its corresponding ground-truth token, without explicitly leveraging the more reliable next-token prediction head to guide the optimization of later prediction heads. This issue becomes more significant as the prediction horizon increases, since predicting more distant future tokens is inherently more difficult.

Although previous studies \cite{samragh2025future} have introduced consistency-based objectives for multi-token prediction, the reliability of the teacher predictions used for consistency supervision has received relatively limited attention. In particular, directly enforcing consistency at low-confidence positions may introduce unreliable supervisory signals. Furthermore, consistency across both output probability distributions and latent representations has not been systematically investigated within the same MTP training framework. These limitations motivate the development of a selective consistency mechanism for improving the optimization of future-token prediction heads.